\documentclass[11pt,letterpaper]{article}
\usepackage[T1]{fontenc}
\usepackage{lmodern,microtype,amsmath,amssymb,amsthm,booktabs,xcolor}
\usepackage[margin=0.94in,top=0.83in,bottom=0.85in,headheight=14pt]{geometry}
\usepackage{fancyhdr}
\usepackage[colorlinks=true,linkcolor=teal!70!black,urlcolor=teal!70!black,citecolor=teal!70!black]{hyperref}
\newtheorem{theorem}{Theorem}
\newtheorem{lemma}[theorem]{Lemma}
\newtheorem{corollary}[theorem]{Corollary}
\newcommand{\C}{\mathbb C}\newcommand{\R}{\mathbb R}
\newcommand{\Z}{\mathbb Z}\newcommand{\A}{\mathcal A}
\newcommand{\rpart}{\operatorname{Re}}\newcommand{\ipart}{\operatorname{Im}}
\newcommand{\dd}{\mathrm d}\newcommand{\ii}{\mathrm i}
\definecolor{ink}{HTML}{16384A}
\pagestyle{fancy}\fancyhf{}\lhead{\small\color{ink}Prime frequencies of differentiated Dirichlet completions}\rhead{\small Research note}\cfoot{\thepage}
\setlength{\emergencystretch}{1.5em}
\hypersetup{pdftitle={Hearing a differential operator and a conductor at one prime},pdfauthor={OpenAI}}
\begin{document}
\thispagestyle{plain}
\noindent{\sffamily\small\bfseries\color{teal!70!black} ANALYTIC NUMBER THEORY\hfill NOVEMBER 2026}
\medskip

\noindent{\LARGE\bfseries\color{ink}Hearing a differential operator and a conductor at one prime\par}
\medskip
\noindent Given the zeros of a derivative, can one tell which derivative was taken?
For a primitive Dirichlet character $\chi$, consider the entire completion
$\Xi_\chi(s)$ and the zeros of $P(\dd/\dd s)\Xi_\chi(s)$, where the monic
polynomial $P$ has purely imaginary roots. We prove that the frequency of a single unramified prime
in these zero ordinates determines $P$. If $\deg P>0$, the same
measurement determines the conductor of $\chi$, even when it is unknown.
The measurable quantity is a rational shift quotient of $P$ to every
order in $1/\log T$. For degree $d$, the coefficient at order $d+1$
completes the reconstruction, and examples show that this last coefficient
can be necessary. As a local consequence, any two distinct primitive-character
and monic-operator pairs have different zero multisets in every sufficiently
high interval of length $C\log\log T$, with a constant depending on the pair.

\section{What one prime measures}
Let $q$ be the conductor, $a\in\{0,1\}$ the parity, and set
\[
 A(s)=(q/\pi)^{(s+a)/2}\Gamma((s+a)/2),\qquad \Xi_\chi(s)=A(s)L(s,\chi).
\]
For the conductor-one character take $a=0$ and include $s(s-1)/2$ in $A$,
so that $\Xi_1=\xi$ is entire. Here the argument of $\Xi_\chi$ is always
$s$, including when it is differentiated. Write $D=\dd/\dd s$, let
$P(X)=\prod_{j=1}^d(X-\ii\alpha_j)$ with $\alpha_j\in\R$, and allow $P=1$.
Put $F=P(D)\Xi_\chi$. Its zeros are $1/2+\ii\gamma$ with $\gamma\in\R$,
as shown in Section~\ref{sec:line}. All sums count multiplicity.

For a fixed prime $p$, set $h=\log p$ and define the observable
\begin{equation}\label{obs}
 S_p(T,H)=\frac{\sqrt{2\pi p}}{H}\sum_\gamma
      \exp\!\left(-\frac{(\gamma-T)^2}{2H^2}+\ii h\gamma\right).
\end{equation}
Negative ordinates are included. For $T\to+\infty$ their contribution
is negligible in the windows below. For a complex character the sign of the
ordinate and the sign of the phase matter. Put
\begin{equation}\label{coord}
 z=\tfrac12\log(T/(2\pi))+\ii\pi/4,\qquad
 c=\tfrac12\log q,\qquad Q(X)=P(X+c).
\end{equation}
\begin{theorem}[Prime-frequency formula]\label{th:trace}
For fixed $\chi,P,p$ there is $\eta_p>0$ such that, uniformly for
$1\le H\le T^{2/3}$ as $T\to+\infty$,
\begin{equation}\label{trace}
 S_p(T,H)=-h\chi(p)\frac{Q(z-h)}{Q(z)}
              +O_{\chi,P,p}\!\left(T^{-1/5}+ (\log T)e^{-\eta_p H^2}\right).
\end{equation}
In particular $H=\log\log T$ gives an error $o((\log T)^{-k})$ for
every fixed $k$. The formula includes $p\mid q$, where $\chi(p)=0$.
\end{theorem}
\begin{theorem}[One-prime reconstruction]\label{th:inverse}
Fix a prime $p\nmid q$. The values of \eqref{obs} with $H=\log\log T$
along any unbounded sequence $T\to\infty$ determine $\chi(p)$ and the
monic $P$. For $d>0$ they also determine $q$. In the expansion in $1/z$,
the constant term determines $\chi(p)$, the first correction gives $d$,
and the second gives $q$ for $d>0$. Terms through $z^{-d-1}$ suffice
for all coefficients of $P$. For each $d\ge1$, terms only through
$z^{-d}$ can agree for distinct operators of degree $d$, even with $q$
and $\chi$ fixed.
\end{theorem}
The theorem uses exact asymptotic data. It imposes no simplicity or
spacing assumption on zeros. The leading coefficient of an arbitrary
nonzero $P$ multiplies the entire function and is invisible to its zeros,
which explains the monic normalization. If $P=1$, all the inverse-log
coefficients beyond the constant vanish, irrespective of the conductor.

\section{Critical-line zeros after differentiation}\label{sec:line}
The arithmetic input is the critical-line theorem for primitive Dirichlet
$L$-functions from \cite{lean}. Classical completion, the Euler product,
Stirling's formula and the Dirichlet zero count are used in their ordinary
forms \cite{dav,bmor}. In particular the zeros of the entire completion
are precisely the nontrivial zeros. The functional equation is
$\Xi_\chi(s)=\varepsilon_\chi\overline{\Xi_\chi(1-\bar s)}$ with
$|\varepsilon_\chi|=1$.
\begin{lemma}\label{interlace}
$F$ has order at most one and every zero lies on $\rpart s=1/2$.
Its number of ordinates in $[u,u+1]$ is $O_{\chi,P}(\log(2+|u|))$.
Moreover $F(s)=(-1)^d\varepsilon_\chi\overline{F(1-\bar s)}$.
\end{lemma}
\begin{proof}
Take a unit square root of $\varepsilon_\chi$ and divide
$\Xi_\chi(1/2+\ii w)$ by it. This is a real entire function $f(w)$
of order at most one, with only real zeros. The usual Dirichlet zero
count gives infinitely many zeros on both tails and the indicated unit
interval bound for $f$. By the genus-one Hadamard product,
\[
 \frac{f'(w)}{f(w)}
   =B+\frac{m_0}{w}
     +\sum_{t\ne0}m_t\left(\frac1{w-t}+\frac1t\right),
 \qquad B\in\R .
\]
Here $t$ runs over distinct nonzero zeros, $m_t$ is their multiplicity and
$\sum_{t\ne0}m_t/t^2<\infty$. The imaginary part of the displayed
expression is $-\ipart(w)\sum_t m_t/|w-t|^2$. Consequently
$(\dd/\dd w+\alpha)f$ has only real zeros for real $\alpha$. At a zero
of order $m$ it retains multiplicity $m-1$. In each intervening gap
$f'/f$ decreases strictly from $+\infty$ to $-\infty$, giving one
additional simple zero. The interval counts on half-open intervals
therefore differ by at most one after each step. Closed unit intervals
are covered by slightly longer half-open intervals. Derivatives preserve
the order bound by Cauchy's estimate, and there are still infinitely
many zeros in both directions. Each $D-\ii\alpha$ corresponds to
$-\ii(\dd/\dd w+\alpha)$. Induction proves the zero assertions. The
last identity follows by conjugation and the chain rule for each factor.
\end{proof}

\section{Proof of the prime-frequency formula}
We give the contour and coefficient computations separately. The choice
of a prime makes the coefficient calculation particularly short.
\subsection{Freezing the gamma factor}
Fix the edge $\sigma=2$. In the Banach algebra $\A_2$ of coefficient
sequences with $\|u\|=\sum |u_n|n^{-2}<\infty$, the product is Dirichlet
convolution. Evaluation at $2+\ii t$ has norm at most one. Write
$g=z+c$ and introduce the two Dirichlet series, with $g$ held fixed,
\[
 U_g(s)=\sum_{n\ge1}\chi(n)P(g-\log n)n^{-s},\qquad
 V_g(s)=-\sum_{n\ge1}\chi(n)(\log n)P(g-\log n)n^{-s}.
\]
The coefficients of $U_g/g^d$ tend in $\A_2$ to those of $L(s,\chi)$
with error $O(1/|g|)$ by polynomial expansion and summability of log
powers. $L^{-1}$ belongs to $\A_2$ by the Euler product. A Neumann
series thus yields $\|U_g^{-1}\|=O(|g|^{-d})$. Since
$\|V_g\|=O(|g|^d)$, the series
\[
 B_g=V_g/U_g=\sum_{n\ge1} b_n(g)n^{-s}
 \quad\text{satisfies}\quad \|B_g\|=O(1).
\]
The coefficient at $1$ is zero. Comparing coefficients in $B_g U_g=V_g$
at the prime $p$ gives the exact identity
\begin{equation}\label{prime}
 b_p(g)=-(\log p)\chi(p)\frac{P(g-\log p)}{P(g)}.
\end{equation}

For $s=2+\ii t$ and $|t-T|\le T^{3/4}$, Stirling's differentiated
formula on this vertical strip gives
\[
  A'/A=g+O(T^{-1/4}),\qquad
  (A'/A)^{(j)}=O_j(T^{-j})\quad(j\ge1).
\]
The error includes the variation of $\log t$ across the interval. The
extra factor in conductor one changes $A'/A$ by $O(T^{-1})$.
By induction using $A^{(k+1)}/A=(A^{(k)}/A)'+(A'/A)(A^{(k)}/A)$,
$A^{(k)}/A$ is $(A'/A)^k$ plus terms containing at least one differentiated
logarithmic derivative. Hence for each fixed $k$,
$A^{(k)}/A=g^k+O_k(T^{-1/4}(1+|g|)^k)$.
In $F/A$ apply $P(D-\log n)$ to $A$ term by term. The corresponding
calculation in $F'/A$ uses $(D-\log n)P(D-\log n)$. The coefficient
errors after replacing $D$ acting on $A$ by $g$ are bounded in $\A_2$
by $O(T^{-1/4}(1+|g|)^{d+1})$ for both expressions, since
$\sum n^{-2}(1+\log n)^{d+1}$ converges. The frozen expressions are
respectively $U_g$ and $gU_g+V_g$. The inverse bound above gives
$|U_g(2+\ii t)|\gg |g|^d$. Division, absorbing fixed powers of
$\log T$ into $T^{1/20}$, proves
\begin{equation}\label{freeze}
 F'/F(2+\ii t)=g+B_g(2+\ii t)+O(T^{-1/5}).
\end{equation}
Using Stirling at each height separately gives
$F'/F(2+\ii t)=O(\log(2+|t|))$ globally. For negative heights the
imaginary term in $g$ changes sign, and compact intervals are bounded by
Lemma~\ref{interlace}.

\subsection{Selecting a single Dirichlet coefficient}
Apply the residue theorem on the edges $\sigma=2,-1$ to $(F'/F)(s)K(s)$,
where
\[
 K(s)=p^s\exp((s-1/2-\ii T)^2/(2H^2)).
\]
Multiply the resulting zero sum by $\sqrt{2\pi}/H$. Its value is exactly
$S_p(T,H)$. On each vertical edge the oriented integral is consequently
the integral in $t$ divided by $\sqrt{2\pi}H$, with the left edge
subtracted. To justify removing the horizontal edges, select heights
tending to $+\infty$ and $-\infty$ that stay $\gg1/\log(2+|t|)$ from
every zero ordinate; this follows from Lemma~\ref{interlace} by deleting
small intervals in each unit interval. The logarithmic derivative there
is $O(\log^2(2+|t|))$ across the rectangle. Indeed, the genus-one sum
at ordinates $|\gamma|\le 2|t|$ is bounded by unit shells about $t$
using the count in the lemma, with the nearest shell bounded separately.
Its regularizing terms sum to $O(\log^2(2+|t|))$. The complementary tail
is $O(|t|\sum_{|\gamma|>2|t|}m_\gamma/\gamma^2)=O(\log(2+|t|))$,
where here $\gamma$ indexes distinct ordinates. The Gaussian then kills
both horizontal integrals. The same count proves absolute convergence
of the sum in \eqref{obs}.

On the right edge substitute \eqref{freeze}; on the left use
$(F'/F)(-1+\ii t)=-\overline{(F'/F)(2+\ii t)}$. Since
$H\le T^{2/3}$ the discarded Gaussian tails $|t-T|>T^{3/4}$ are
smaller than any inverse power of $T$, by the global edge bound and
the uniform norm of $B_g$. The integrated freezing error is $O(T^{-1/5})$.
For a term with frequency $v$ on an edge $\sigma$ the normalized
Gaussian integral, stripping constant factors from that term, is exactly
\[
 \frac1{\sqrt{2\pi}H}\int_{\R}e^{\ii vt}
   e^{(\sigma-1/2+\ii(t-T))^2/(2H^2)}\,\dd t
   =e^{\ii vT}e^{-H^2v^2/2-(\sigma-1/2)v}.
\]
On the right $v=\log(p/n)$ and on the left $v=\log(pn)$ for the
conjugate series. The constant terms $g,\bar g$ have $v=h$. There is
one resonance, the right-edge term $n=p$; its edge factor $p^2p^{-2}$ is
one. All other frequencies have modulus at least $\delta=\log(1+1/p)$.
For either edge and $H\ge1$ the Gaussian multiplier is bounded by
$C\exp(-H^2 v^2/4)$, by completing the square in the remaining exponent.
The total discarded contribution is $O((\log T)e^{-\delta^2H^2/4})$
using the norm of $B_g$. Inserting \eqref{prime} proves
Theorem~\ref{th:trace}, for example with $\eta_p=\delta^2/4$.

\section{Inverting the shift quotient}
Divide the measured expansion by its limit $-h\chi(p)$, which is
nonzero at every unramified prime. Write the result as
$1+\sum_{j\ge1}r_jz^{-j}$. Successive coefficients are uniquely recovered
by subtracting known terms, multiplying by $z^j$ and passing to the
limit. Theorem~\ref{th:trace} permits this along any unbounded sequence.
For $Q=X^d+a_1X^{d-1}+\cdots+a_d$, direct expansion gives
\begin{equation}\label{jet}
 r_1=-dh,\qquad r_2=\tbinom d2 h^2+h a_1\quad(d\ge1).
\end{equation}
The roots of $Q$ are $-c+\ii\alpha_j$. In particular
$\rpart a_1=dc$, whence an explicit conductor formula is
\begin{equation}\label{cond}
 q=\exp\!\left(\frac{2}{dh}\rpart\left[r_2-\tbinom d2 h^2\right]\right).
\end{equation}

For the higher coefficients, suppose $Q_1,Q_2$ are monic of the same
degree and their first difference is $\Delta a_j$ in position $j$.
Then $Q_1/Q_2=1+\Delta a_j z^{-j}+O(z^{-j-1})$. Taking a logarithm
near infinity and subtracting its expansions at $z-h$ and $z$ shows
\begin{equation}\label{tri}
 \frac{Q_1(z-h)}{Q_1(z)}-\frac{Q_2(z-h)}{Q_2(z)}
       =jh\,\Delta a_j\,z^{-j-1}+O(z^{-j-2}).
\end{equation}
Indeed the difference of each subsequent Laurent monomial gains one
inverse power. Exponentiating preserves the first nonzero difference
since both shift quotients tend to one. To compute $a_j$, form a trial
polynomial with the previously recovered coefficients and all remaining
coefficients zero. Subtract its coefficient at $z^{-j-1}$ from $r_{j+1}$
and divide by $jh$. This recovers $Q$ and hence $P(X)=Q(X-c)$.

To prove sharpness, let $f(x)=\prod_{j=1}^d(x-j)$. For sufficiently small
nonzero real $t$, $f+t$ also has $d$ simple real roots, by preserving the
signs at both ends of disjoint intervals about the old roots. Set
$P_1(X)=\ii^d f(X/\ii)$ and $P_2(X)=\ii^d(f(X/\ii)+t)$.
The polynomial difference is constant, and translation by $c$ preserves it.
Equation~\eqref{tri} gives equality through $z^{-d}$ and inequality
at $z^{-d-1}$. This proves Theorem~\ref{th:inverse}.

\paragraph{A quadratic character example.}
Take the quadratic character modulo $5$, $p=2$ and
$P(X)=(X-2\ii)(X+\ii)=X^2-\ii X+2$. Then $h=\log 2$,
$\chi(2)=-1$, $a_1=2c-\ii$, $a_2=c^2-\ii c+2$. The limit of $S_2$
is $+h$. Its normalized first correction $-2h$ gives degree two, and
$r_2=h^2+(2c-\ii)h$ gives $q=5$ by \eqref{cond}. The third correction
$(2a_2-a_1^2)h-a_1h^2$ gives $a_2$ and finishes the reconstruction.
The script \texttt{check\_algebra.py} performs these operations exactly.

\section{How local is the information?}
The Gaussian in \eqref{obs} uses every ordinate. For any prescribed
finite order of accuracy the observations can be confined to an interval.
\begin{corollary}\label{local}
For every fixed $k\ge0$ there exist $A,B>0$ such that the sum
\eqref{obs} with $H=A\sqrt{\log\log T}$, restricted to
$|\gamma-T|\le B\log\log T$, equals the rational main term in \eqref{trace}
up to $o((\log T)^{-k})$.
For two distinct pairs $(\chi_1,P_1),(\chi_2,P_2)$ as in
Theorem~\ref{th:trace}, there is $C>0$ such that their zero ordinates
differ as multisets in $[T-C\log\log T,T+C\log\log T]$ for every
sufficiently large $T$. In $[X,2X]$ there are consequently
$\gg X/\log\log X$ unmatched zero occurrences, counted with multiplicity.
The constants here may depend on both pairs.
\end{corollary}
\begin{proof}
Take $\eta_p A^2>k+2$. Write $L=\log\log T$ and $R=BL$.
By the local zero count and summation in unit intervals, for $1\le H\le
T^{2/3}$ the sum of $e^{-(\gamma-T)^2/(4H^2)}$ is $O(H\log T)$:
for $|\gamma-T|\le T$ use the bound $O(\log T)$ per interval, and
beyond $T$ use Gaussian decay against $\log(2+|\gamma|)$.
Pulling out $e^{-R^2/(4H^2)}$ from the original Gaussian bounds the
normalized omitted tail by $O((\log T)e^{-B^2 L/(4A^2)})$.
Take $B^2/(4A^2)>k+2$.

Distinct primitive characters have different values at some prime.
Otherwise complete multiplicativity makes the underlying arithmetic
functions equal, with the same primitive conductor and character.
At such a prime the constant terms differ. For equal characters choose
any unramified prime. Different monic polynomials then give different
quotients at a finite order $k$ by \eqref{jet} and \eqref{tri}. In each
case the trace difference has a first term of size
$\asymp(\log T)^{-k}$, including $k=0$ for different characters.
A common multiset on the truncated interval would give identical
truncated traces. This contradicts the first assertion. Finally pack
$\gg X/\log\log X$ disjoint intervals of the stated size inside
$[X,2X]$. Each contains a discrepancy of at least one in integer
multiplicity.
\end{proof}
One prime need not distinguish characters that agree there. For instance,
distinct primitive characters of the same conductor with identical
values at $p$ give the same rational term for a fixed $P$. Across all primes
the leading terms determine the character. The short-interval conclusion
allows the discrepant ordinate to be a common zero of different multiplicity.

\paragraph{Finite-height estimation.}
Suppose $d$ is known and put $m=d+1$. At centers $T^j$ for
$j=1,\ldots,m+1$, use the truncated sums above with sufficiently large
accuracy, say $k=m+1$. At each center the main term through degree $m$
in $1/z$ has residual $O((\log T)^{-m-1})$. Put $L_0=\log T$. The
nodes $L_0/z(T^j)$ tend to $2/j$; the associated fixed-size Vandermonde
matrices have uniformly bounded inverses. Interpolating the $m+1$ sums
therefore estimates the coefficient of $z^{-k'}$ with error
$O(L_0^{k'-m-1})$ for $0\le k'\le m$, after multiplying the
interpolated coefficients by $L_0^{k'}$. At an unramified prime the
estimated constant is bounded away from zero, so division and the
triangular inversion above yield consistent coefficient estimates.
For $d>0$ the estimated conductor eventually rounds to the correct
integer. This is an asymptotic assertion; constants depend on the
fixed data and no universal numerical height threshold is claimed.

\section{A version using full complex zeros}
There is also an independent route when the observations give both
coordinates of each zero. It clarifies the precise role of Dirichlet GRH.
\begin{theorem}\label{complex}
Without using the critical-line theorem, form the sum
\[
 \widetilde S_p(T,H)=\frac{\sqrt{2\pi}}{H}
 \sum_{F(\rho)=0}p^\rho\exp\big((\rho-1/2-\ii T)^2/(2H^2)\big)
\]
over full complex zeros with multiplicity. Formula~\eqref{trace} and
the reconstruction of Theorem~\ref{th:inverse} hold for this sum.
Corollary~\ref{local} holds for complex zero multisets with ordinates
in the indicated intervals.
\end{theorem}
\begin{proof}
We supply the two replacements for Lemma~\ref{interlace}. First the
zeros of $\Xi_\chi$ lie in $0\le\rpart s\le1$, by the Euler product and
functional equation. In the $w$ coordinate of that lemma this is
$|\ipart w|\le1/2$. For a real entire function of order at most one
with its zeros $v$ in that strip, group the genus-one product in conjugate
pairs. The regularizing constants have zero imaginary part when paired.
Consequently
\[
 \ipart(f'/f)(w)=-\sum_v\frac{\ipart w-\ipart v}{|w-v|^2},
\]
including the contribution at zero. This series is absolutely convergent
at points off the zero set. If there is at least one zero, the sum has
strict sign in either exterior half-plane, so $f'+\alpha f$ cannot vanish
there. A zero-free intermediate function of order at most one is an
exponential of a linear polynomial and its transform is either zero-free
or identically zero. The latter is excluded at each step by the right-edge
estimate $|P(D)\Xi_\chi(2+\ii t)/A(2+\ii t)|\gg(\log |t|)^d$ for the
corresponding partial product, established in the freezing calculation
by $U_g^{-1}$ without any zero-location assumption. This proves strip
confinement through all factors.

Second, on a disk of radius $5$ centered at $2+\ii u$, for large $|u|$,
the supremum of $|F|$ is $\ll e^{-\pi|u|/4}(2+|u|)^M$ for fixed $M$.
This follows from the polynomial vertical-strip growth of $L$, Stirling
on a slightly larger disk, and Cauchy's derivative estimate. At the
center, the same right-edge lower bound and Stirling give
$|F|\gg e^{-\pi|u|/4}(2+|u|)^{-M}$ after increasing $M$.
Jensen's formula between radii $3$ and $5$ gives $O(\log(2+|u|))$
zeros in the disk of radius $3$, covering the strip segment at heights
$u$ through $u+1$. Compact heights cause no difficulty. With strip
confinement and this count, the horizontal contour bounds follow by the
same unit-shell estimate. Reflection still gives the left edge.
The residues are now the summands of $\widetilde S_p$. Their absolute
values before normalization are bounded by
$p\exp(1/(8H^2))e^{-(\ipart\rho-T)^2/(2H^2)}$, so the truncation
proof also applies. All remaining calculations concern the right edge
or the polynomial and carry over verbatim.
\end{proof}

\section{Comparison with sources and numerical checks}
\paragraph{The formal input and ProveIt.}
At commit \texttt{72fead102f66}, the file \texttt{grh/src/Main.lean}
exports \texttt{riemannHypothesis}, which proves the proposition
\texttt{RiemannHypothesis} from Mathlib. Its primitive-character theorem
is \texttt{LFunction\_\allowbreak eq\_\allowbreak zero\_\allowbreak re\_\allowbreak of\_\allowbreak isPrimitive}.
The file also exports half-plane nonvanishing for all Dirichlet characters
\cite{lean}.
One bridge in the development uses a one-sided estimate for smoothed
twisted prime sums to exclude off-line zeros. Here the direction of
inquiry is reversed: a smoothed sum over the zeros of a differentiated
completion reveals a prime coefficient. We use the exported zero-location
result, the functional equation, and the Euler product. The Lean
development reports its own kernel and independent checks. This note is
a mathematical proof on paper; the upstream build was not replayed here.

The requested ProveIt index at \texttt{2c45d198070f} describes 61 research
packets and states their archival status \cite{proveit}. Among the additional
incoming archives, \texttt{grh\_spectral\_rigidity.zip} studies weighted
ordinary zero divisors, and \texttt{prime\_zero\_reconstruction.zip}
studies the recovery of zero data from a signal assembled from primes.
In the present inverse problem the unknown is the operator applied to the
completed function before its zeros are observed. The requested Wikipedia
article discusses the November 2026 formal proof and the classical
relation between prime oscillations and zeta zeros \cite{wiki}.

\paragraph{Published precedents.}
Farmer and Gonek's analysis of $\xi'$ already contains the degree-one
prime correction \cite[Lemma 3.1]{fg}. Malik and Roy give Landau-type
formulas for $\xi^{(m)}$ with uniformity in $m$ \cite[Theorem 1.3]{mr}.
Srinivas proves a $\log\log T$ short-interval distinction for ordinary
Selberg-class functions under the hypotheses in his theorem \cite{sr};
Banks gives a conductor-dependent comparison for primitive characters
\cite{banks}. These precedents concern central ingredients here.
What we have not located in them or the inspected ProveIt work is the
recovery of an unknown polynomial and unknown conductor from one
unramified frequency, with the sharp finite order of reconstruction,
or the associated short-interval comparison for two different operators
on the same completion. Priority for these statements remains open to
further bibliographic review.

\paragraph{Numerical example.}
For $P(X)=X$ in conductor one, \eqref{trace} predicts
$-h(1-h/z)$. Two direct sums at $p=3$ give the following values.
\begin{center}\small
\begin{tabular}{rrrrl}
\toprule
$T$ & $H$ & $\xi'$ ordinates & \multicolumn{1}{c}{Observed $S_3$} &
\multicolumn{1}{c}{Predicted} \\
\midrule
1500&45&626& $-0.691166-0.116778\ii$ & $-0.691271-0.116860\ii$ \\
6000&90&1572& $-0.764319-0.076509\ii$ & $-0.764334-0.076524\ii$ \\
\bottomrule
\end{tabular}
\end{center}
The corresponding sums for $P=1$ are $-1.098612$ to six decimal places
in both windows, as predicted by $-\log 3$. The script
\texttt{zero\_example.py} uses Euler--Maclaurin on the critical line,
checks the number of sign-change brackets against an independent
\texttt{mpmath} zero count, solves for the critical points between
successive zeros of $\xi$, and retains the ordinates within $8H$.
Its output is a floating-point check of the sign, normalization and
first-derivative correction. The algebra script separately checks the
triangular inverse and the last-coefficient examples.

\paragraph{Further questions.}
Can the length $C\log\log T$ be replaced by a bounded length for
two different operators on the same completion? How small can the
constants in the finite-height inverse problem be made uniformly in
the degree and conductor? For higher-degree $L$-functions, the
archimedean logarithmic derivative introduces several gamma parameters.
Which of them can be separated at a single prime frequency, and which
require accuracy finer than every inverse power of $\log T$?

\begin{thebibliography}{99}\small
\bibitem{lean} L. Alp\"oge, \emph{GRH for Dirichlet $L$-functions, in Lean 4},
\href{https://github.com/lsernau/grh-lean/tree/72fead102f663d5c7a407964a13dc0cf9d977207}{pinned repository},
\href{https://github.com/lsernau/grh-lean/blob/72fead102f663d5c7a407964a13dc0cf9d977207/grh/src/Main.lean}{theorem exports}.
\bibitem{proveit} V. Reshetnikov, \emph{ProveIt, Generalized Riemann Hypothesis},
\href{https://github.com/VladimirReshetnikov/ProveIt/tree/2c45d198070f93209d6f0faa2d387ef0a2c4e72c/Analysis/GeneralizedRiemannHypothesis}{directory and provenance index}.
\bibitem{wiki} Wikipedia contributors, \emph{Riemann hypothesis},
\href{https://en.wikipedia.org/wiki/Riemann_hypothesis}{article}, consulted November 2026.
\bibitem{dav} H. Davenport, \emph{Multiplicative Number Theory}, 3rd ed.,
revised by H. L. Montgomery, Springer, 2000, chapters 9--16.
\bibitem{bmor} M. A. Bennett, G. Martin, K. O'Bryant and A. Rechnitzer,
\href{https://doi.org/10.1090/mcom/3599}{Counting zeros of Dirichlet $L$-functions},
Math. Comp. 90 (2021), 1455--1482.
\bibitem{fg} D. W. Farmer and S. M. Gonek,
\href{https://arxiv.org/abs/0803.0425}{Pair correlation of the zeros of the derivative of the Riemann $\xi$-function}, preprint (2008). Published with Y. Lee in \href{https://doi.org/10.1112/jlms/jdu026}{J. London Math. Soc. 90 (2014), 241--269}.
\bibitem{mr} A. Malik and A. Roy,
\href{https://doi.org/10.1515/forum-2018-0081}{On the distribution of zeros of derivatives of the Riemann $\xi$-function}, Forum Math. 32 (2020), 1--22.
\bibitem{sr} K. Srinivas,
\href{https://www.imsc.res.in/~srini/publication/recent/selberg.pdf}{Distinct zeros of functions in the Selberg class}, Acta Arith. 103 (2002), 201--207.
\bibitem{banks} W. D. Banks,
\href{https://arxiv.org/abs/2302.07073}{Comparing zeros of distinct Dirichlet $L$-functions}, arXiv:2302.07073, version 2 (2024).
\end{thebibliography}
\end{document}
